% =====================================================================
% INTERNAL FIRST DRAFT -- NOT FOR DISTRIBUTION
% qec-ai-search, paper/draft.tex
% Every quantitative claim carries a tag [Cn]; see paper/claims_map.md for
% the mapping  Cn -> (research/STATE.md "confirmed" item, audit file + section,
% results path).  Set \showclaimtagsfalse to hide the tags.
% Anything that could not be traced to the project files is marked [UNVERIFIED].
% Compilation has NOT been tested (see paper/README.md).
% =====================================================================
\documentclass[aps,prl,twocolumn,superscriptaddress,nofootinbib]{revtex4-2}

\usepackage{amsmath,amssymb}
\usepackage{booktabs}
\usepackage{xcolor}
\usepackage[colorlinks=true,linkcolor=blue,citecolor=blue,urlcolor=blue]{hyperref}

% ---- internal-draft macros ------------------------------------------------
\newif\ifshowclaimtags \showclaimtagstrue
\newcommand{\cl}[1]{\ifshowclaimtags\textsuperscript{\textcolor{blue}{\scriptsize[C#1]}}\fi}
\newcommand{\todo}[1]{\textcolor{red}{[TODO: #1]}}
\newcommand{\unverified}{\textcolor{red}{[UNVERIFIED]}}

\newcommand{\dcirc}{d_{\mathrm{circ}}}
\newcommand{\dext}{d_{\mathrm{ext}}}
\newcommand{\dt}{\tilde d}
\newcommand{\code}[1]{[[#1]]}

\begin{document}

\title{Circuit distances of bivariate-bicycle syndrome circuits must be proven, not inferred}

\author{\todo{author list -- PI decision}}
\affiliation{\todo{affiliations}}

\date{Internal draft, \today}

\begin{abstract}
The circuit-level distance $\dcirc$ of a syndrome-extraction circuit sets the
low-noise scaling of its logical error rate, yet for bivariate-bicycle (BB) codes
and their morphing variants it has almost always been reported as heuristic upper
bounds, conjectures, or Monte-Carlo beliefs. We compute exact circuit distances,
with complete lower-bound certificates, for a set of BB and morphing circuits
under circuit-level noise with independent single-operation faults. (i) The
Shaw--Terhal morphing circuit derived from the $\code{288,12,18}$ code, acting on
a $\code{144,12,12}$ end-cycle code, has $\dcirc=12$ for any number of rounds:
a fully distance-preserving circuit without dedicated ancillas. (ii) Of all 80
two-round morphing circuits of the same construction, 72 preserve the end-cycle
distance and 8 do not; for $\code{54,8,8}$ the circuit distance depends on a term
ordering that the construction leaves free, contrary to the order-independence
stated for these circuits. (iii) The Strikis--Browne--Beverland gross-code
circuit has $\dcirc=11$, proving their conjecture; within an explicit family of
period-$\le 8$ retimings that keep its CNOT orders, no circuit reaches 12, and we
give scoped bulk obstructions beyond it. (iv) A depth-8 single-ancilla circuit for
$\code{108,8,10}$ has $\dcirc=10$ (previously $\le 8$), and further BB circuits
have proven distances 6, 9, 10 and 9. The methods (extended codes with
exhaustive symmetric search, an exact lift test, a computationally verified
time-folding reduction, and frame-model certificates) are general.
\end{abstract}

\maketitle

% ---------------------------------------------------------------------------
\paragraph{Introduction.}
Bivariate-bicycle (BB) codes~\cite{Bravyi2024} are leading candidates for
low-overhead quantum memory. What matters in practice, however, is not the code
distance $d$ but the circuit-level distance $\dcirc$: the minimum number of
circuit faults that flip a logical observable without triggering any detector.
At low physical error rate $p$ the logical error rate of a circuit is
governed by $\dcirc$ (it scales as $p^{\lceil \dcirc/2\rceil}$ for a
minimum-distance decoder), so $\dcirc$ rather than $d$ is the figure of merit of
a syndrome circuit.

For BB circuits, essentially every reported $\dcirc$ is an upper bound obtained by
randomized minimum-weight search (e.g.\ BP-OSD)~\cite{Bravyi2024,ShawTerhal2025,TourDeGross},
and the most-cited values are conjectures: IBM conjectured its
$\code{72,12,6}$ circuit to be distance preserving~\cite{Bravyi2024};
Strikis, Browne and Beverland conjectured their gross-code circuit to have
$\dcirc=11$~\cite{Strikis2026}; and a morphing circuit for a gross-sized code is
``believed'' to reach 12 on the basis of Monte-Carlo sampling~\cite{Wills2026}.
Upper bounds are cheap and lower bounds are hard, so it is tempting to treat
an unrefuted upper bound as the answer, or to transfer a property that holds for
the underlying \emph{code} to the \emph{circuit}.

This Letter shows that both shortcuts can fail, and that exact answers are
within reach. We prove exact $\dcirc$ for a set of BB and morphing circuits
(Table~\ref{tab:main}), find circuits where the inferred answer is wrong, and
delimit precisely what can and cannot be concluded about reaching $\dcirc=12$
for the gross code with single-ancilla circuits.

% ---------------------------------------------------------------------------
\paragraph{Setting.}
We consider memory experiments: data initialisation, $R$ rounds of syndrome
extraction, and transversal data readout, in both the $Z$ and $X$ basis.
Faults are single-operation Pauli components of standard circuit-level
depolarizing noise (SD6~\cite{Gidney2021honeycomb}: two-qubit depolarization after each CNOT, flips after
resets and before measurements, single-qubit depolarization on idle qubits),
counted as distinct detector-error-model (DEM) mechanisms~\cite{Gidney2021stim}.
$\dcirc$ is the minimum number of such faults whose combined effect triggers no
detector and flips a logical observable; we report it per basis, and in every case
below the two bases agree. All lower bounds below also hold for any noise model
whose faults are a subset of these (e.g.\ no idle noise). Correlated
multi-operation faults, leakage, and CZ-based compilations are outside scope.
Upper bounds are always realised by explicit witnesses injected with probability
one into the noiseless \texttt{stim} circuit (zero detectors fired, an
observable flipped).

\paragraph{Methods in brief.}
\emph{(a) Extended code and exhaustive search.} For single-ancilla CNOT circuits,
Strikis \emph{et al.}~\cite{Strikis2026} append to the parity-check matrix one
column per ``residual'' (a suffix of a check's CNOT order), giving an extended code
with distance $\dext$; their Prop.~3 states $\dcirc\ge\dext$ for any such circuit.
We verify the key step of Prop.~3 directly on each circuit to which we apply it (every single
fault propagates to zero or exactly one extended-code column)\cl{27} and decide
$\dext\ge w$ by a complete depth-first search (DFS) that branches on an
unsatisfied syndrome row with sibling exclusion and splits cases by the orbits of
the circuit's translation symmetry. The DFS was validated against unpruned
meet-in-the-middle enumeration and rerun without symmetry where affordable
(SM~\ref{sm:dfs}).
\emph{(b) Exact lift.} When $\dext<\dcirc$ (interleaved circuits), a weight-$\dext$
undetected circuit error must map to $\dext$ distinct nonzero extended-code
columns; enumerating all minimum-weight extended logicals and testing whether
each lifts to an undetected circuit error decides $\dcirc$ exactly
(SM~\ref{sm:lift}).
\emph{(c) Time folding for morphing circuits.} Shaw and Terhal's Prop.~1~\cite{ShawTerhal2026}
reduces a two-round morphing circuit with any number of rounds to a
single-round instance. Rather than relying on its per-fault time-reversal argument,
we verify by linear algebra the basis-free condition the fold actually needs
(SM~\ref{sm:fold})\cl{3}.
\emph{(d) Frame-model certificates.} For retimed gross circuits we use a bulk
Pauli-frame model whose feasibility implies a genuine weight-11 undetected error,
and certify infeasibility of blocking all such errors with CP-SAT and a
lazy-clause (CEGAR) SAT encoding (SM~\ref{sm:gross12}).

% ---------------------------------------------------------------------------
\begin{table}[t]
\caption{Proven circuit distances. ``Prior'' is the best value previously reported
for the same circuit or code. All entries: SD6 fault set, both memory bases,
any number of rounds $R$ unless noted. Tags refer to \texttt{claims\_map.md}.}
\label{tab:main}
\begin{ruledtabular}
\begin{tabular}{llcc}
code / circuit & prior & $\dcirc$ & tag\\
\hline
$\code{288,12,18}\to\code{144,12,12}$ morphing ($f_x$) & $\le 12$~\cite{ShawTerhal2025} & $\mathbf{12}$ & \cl{1}\\
80 two-round morphing circuits & $\le$ bounds~\cite{ShawTerhal2025} & 72 $=\dt$, 8 $<\dt$ & \cl{5}\\
gross $\code{144,12,12}$, Fig.~12 of~\cite{Strikis2026} & 11 (conj.) & $\mathbf{11}$ & \cl{10}\\
$\code{108,8,10}$, $y$-parity two-colour & $\le 8$~\cite{Bravyi2024} & $\mathbf{10}$ & \cl{19}\\
$\code{72,12,6}$, IBM circuit~\cite{Bravyi2024} & 6 (conj.)$^{a}$ & $6$ & \cl{21}\\
$\code{72,12,6}$, two-colour LRC & -- & $6$ & \cl{22}\\
coprime $\code{126,12,10}$, uniform LRC & $\le 9$~\cite{WangMueller} & $9$ & \cl{22}\\
$\code{120,8,12}$ $(\ell,m)=(20,3)$, uniform LRC & -- & $10$ & \cl{22}\\
$\code{124,10,10}$ $(31,2)$, uniform LRC & -- & $9$ & \cl{22}\\
\end{tabular}
\end{ruledtabular}
\raggedright\footnotesize
$^{a}$An exact value 6 for one basis, $R=6$, and a CZ/SI1000 compilation is
contained in the public data of Ref.~\cite{Webster2026}; see text.
\end{table}

% ---------------------------------------------------------------------------
\paragraph{A fully distance-preserving morphing circuit.}
Morphing circuits~\cite{ShawTerhal2025} have no dedicated ancillas: each round
contracts half of the stabilisers of a ``mid-cycle'' code onto single qubits,
measures them, and re-expands, so that between rounds the state is encoded in an
``end-cycle'' code. Applied to the $\code{288,12,18}$ BB code with the
homomorphism $f_x$, this yields a 288-qubit, degree-5 circuit whose end-cycle code
has parameters $\code{144,12,12}$ and weight-9 checks. Ref.~\cite{ShawTerhal2025}
reports $\dt_{\rm circ}\le 12$ from BP-OSD; Ref.~\cite{Wills2026} (App.~C)
describes a morphing idling circuit ``for the gross code'' that is believed to have
circuit distance 12, supported by Monte-Carlo sampling.

We prove that this circuit (homomorphism $f_x$, layer order of the public
implementation) has $\dcirc=12$ for every number of rounds, in both bases, with
noisy or ideal initialisation and readout\cl{1}. The single-round instance has
144 nontrivial detector rows and 864 projected fault columns with a
72-element translation symmetry; an exhaustive DFS at weight $\le 11$ finds no
undetected logical in either basis, reproduced independently on a second
construction with a different detector basis and search tree
($\approx 2\times10^{11}$ nodes per basis). Twelve-fault witnesses (end-cycle
logical operators, the end-cycle code having distance exactly 12) give the
matching upper bound. The fold condition required to pass from one round to
any $T$ was verified for $T=2,\dots,5$, including the noisy final round, and its
local form is periodic\cl{3}.

Two points of scope matter. First, the end-cycle code is \emph{not} the gross code:
its minimum stabiliser weight is exactly 9, versus 6 for the gross code, so the two
are inequivalent under qubit permutations and local Cliffords\cl{2}. Whether the
circuit of Ref.~\cite{Wills2026} is this one is not stated there and we have not
confirmed it. Second, the statement is for this specific circuit; other variants
are covered below. To our knowledge this is the first rigorous lower bound on
$\dcirc$ for any morphing circuit\cl{4}.

% ---------------------------------------------------------------------------
\paragraph{All 80 two-round morphing circuits.}
The construction of Ref.~\cite{ShawTerhal2025} has more freedom than its tables
suggest. For each of the ten admissible (mother code, homomorphism) pairs it
offers the standard contraction of its Table~IV and the reversed-second-round
variant of its Table~VIII, and within each the order of the terms $(a_2,a_3)$ and
$(b_2,b_3)$ is a writing convention. This gives 80 distinct circuits, which we
reconstructed gate by gate from the paper (Tables IV, VIII, IX) and matched
against an independent implementation\cl{5}. We computed the exact $\dcirc$ of
all 80 in both bases: 72 have $\dcirc=\dt$ (the end-cycle distance) and 8 have
$\dcirc<\dt$ (Table~\ref{tab:morph})\cl{5}.

Three findings stand out. (1)~For $\code{108,8,10}\to\code{54,8,8}$, the four
legal term orders of Table~IV give $\dcirc=6,6,7,7$: the two orders using
$b_3=x^2$ in the first round give 6, the two using $b_3=x$ (including the public
implementation's order) give 7\cl{6}. The caption of Table~IX in
Ref.~\cite{ShawTerhal2025} states that the Table~IV circuits are independent of
the order of the $A$ and $B$ terms; this holds for the end-cycle codes, but not for
the circuit distance\cl{6}. The reported bound $\le 7$ coincides with the exact
value for the public implementation's order. (2)~For $f_x$ applied to the gross
code with the Table~VIII construction, the end-cycle code is $\code{72,12,8}$ but
$\dcirc=6$; Ref.~\cite{ShawTerhal2025} (Table~IX) already reports $\le 6$ for this
circuit, so our contribution is to show that this bound is tight\cl{7}.
Likewise, its $\code{72,12,7}$ circuit ($f_y$) has $\dcirc=7$, also tight\cl{7}.
(3)~The sub-$\dt$ errors are not boundary artefacts: in every one of the 8
circuits and both bases there is a minimum-weight undetected error consisting only
of CNOT faults within a single interior half-cycle, in a model without idle noise\cl{8}.
All 24 circuits derived from $\code{288,12,18}$ preserve their end-cycle distance
(12 or 10); three of these 24 were reproduced by a fully independent pipeline
(SM~\ref{sm:morph80})\cl{9}.

\begin{table}[t]
\caption{Exact $\dcirc$ of the 80 two-round morphing circuits of
Ref.~\cite{ShawTerhal2025} (SD6, both bases equal). $\dt$: end-cycle distance.
Variants: Table~IV (``std'') or Table~VIII (``rev''), with the four
$(a_2,a_3)/(b_2,b_3)$ orders.}
\label{tab:morph}
\begin{ruledtabular}
\begin{tabular}{llccc}
mother & end-cycle (variants) & \# & $\dcirc$ & $=\dt$?\\
\hline
$\code{72,12,6}$ & $\code{36,12,3}$ & 18 & 3 & yes\\
 & $\code{36,12,4}$ (rev) & 6 & 4 & yes\\
$\code{108,8,10}$ & $\code{54,8,8}$ std, $b_3=x^2$ & 2 & 6 & \textbf{no}\\
 & $\code{54,8,8}$ std, $b_3=x$ & 2 & 7 & \textbf{no}\\
 & $\code{54,8,8}$ rev & 2 & 7 & \textbf{no}\\
 & $\code{54,8,7}$ rev & 2 & 7 & yes\\
$\code{144,12,12}$ & $\code{72,12,6}$ & 18 & 6 & yes\\
 & $\code{72,12,5}$ (rev) & 2 & 5 & yes\\
 & $\code{72,12,7}$ (rev, $f_y$) & 2 & 7 & yes\\
 & $\code{72,12,8}$ (rev, $f_x$) & 2 & 6 & \textbf{no}\\
$\code{288,12,18}$ & $\code{144,12,12}$ & 16 & 12 & yes\\
 & $\code{144,12,10}$ (rev) & 8 & 10 & yes\\
\end{tabular}
\end{ruledtabular}
\end{table}

% ---------------------------------------------------------------------------
\paragraph{The gross code: 11 is proven; 12 is obstructed in scoped families.}
Strikis \emph{et al.} proved that no non-interleaved single-ancilla circuit for the
gross code reaches $\dcirc=12$ (their Prop.~4) and presented a depth-8
three-colour left-right circuit (their Fig.~12) conjectured to have $\dcirc=11$~\cite{Strikis2026}.
We prove the conjecture: $\dcirc=11$ for any $R$, both bases\cl{10}. Its extended code
has 360 merged columns per side and a 24-element order-preserving translation
symmetry (15 orbits); a complete DFS at weight 10 finds no logical in either side,
reproduced on an independently constructed lattice-coordinate instance and under a
second, smaller symmetry group, while weight 11 is attained\cl{11}. To our
knowledge this is the first rigorous lower bound on $\dcirc$ for any gross-code
syndrome circuit, and by Prop.~4 of Ref.~\cite{Strikis2026} it is optimal among
non-interleaved circuits\cl{12}.

Can interleaving, i.e.\ retiming the same per-check CNOT orders, recover 12? Each
side has exactly 4224 weight-11 extended-code logicals (176 free orbits under the
24-element group), and every one of them lifts to an undetected circuit error of
Fig.~12\cl{13}. A circuit with $\dcirc=12$ must prevent all of them from lifting.
For the circuit family used in our benchmark---periodic single-ancilla schedules
with period $P\le 8$, one reset, six CNOTs at distinct offsets and one measurement
per check per period---the question is decided exactly. Keeping the CNOT orders of
Fig.~12 forces $P=8$ (weight-6 checks need $P\ge 8$) and leaves one start time per
check; a complete backtracking enumeration finds 32 conflict-free assignments,
i.e.\ 4 up to global translation, all colour-uniform, deterministic and
non-interleaved, Fig.~12 among them\cl{14}. Each of the four has
injection-verified weight-11 undetected errors for every $R=1,\dots,6$ in both
bases\cl{15}. Hence no retiming of Fig.~12 within this family exceeds 11; this
statement is exact, includes the time boundaries, and does not use the frame
model.

Beyond this family we can only make bulk statements (sufficiently many rounds),
via a frame model that agreed with the exact lift checker in the bulk on every
circuit tested, but fails at the initial time boundary\cl{18}. In the bulk, every legal deterministic circuit that
keeps the Fig.~12 CNOT orders has a weight-11 undetected error consisting of ten
data faults and one hook fault, for: translation-symmetric non-wrapping frame
circuits of any period; translation-symmetric pipelined circuits with
$P=8,\dots,11$ (and $P=12$ by a single encoding); pipelined circuits invariant
under any of several translation subgroups at $P=9,10$; and per-check-free
non-wrapping frame circuits with $P\le 16$ (one encoding, reproduced but not
independently re-encoded)\cl{16}. Open are all pipelined families with
$P\ge 9$ not listed (in particular per-check-free circuits, order-2 subgroups,
and the order-4 subgroup $\langle(3,0)\rangle$ at $P=9$), symmetric pipelined
circuits with $P\ge 13$, per-check-free non-wrapping circuits with $P>16$, and any
circuit that changes the CNOT orders\cl{17}. In a classical control,
$\approx1.1\times10^5$ random-walk evaluations of legal deterministic per-check
$P=9$ circuits never blocked a single one of the 720 single-hook logicals in the
frame model (SM~\ref{sm:gross12}); this is evidence, not proof.

% ---------------------------------------------------------------------------
\paragraph{Distance-preserving circuits for other BB codes.}
For $\code{108,8,10}$ the IBM circuit has $\dcirc\le 8$~\cite{Bravyi2024,ShawTerhal2025},
and while the single-check bound for partitionable non-interleaved circuits does
not exclude 10 for this code, uniform depth-8 tilings are limited to
$\le 9$\cl{24}. A two-colour left-right circuit that
reverses the CNOT order within one term block on checks of one $y$-parity (the
$X$ side on odd rows, the $Z$ side on even rows; depth 8, one ancilla per check,
all checks every round) has $\dcirc=10=d$\cl{19}. To our knowledge this is the
first $\dcirc=10$ circuit for this code and the first rigorous circuit distance
for a BB code with $d\ge10$; rigorous circuit distances exist for other qLDPC
families~\cite{ManesClaes,Liang2026,SunCai2026}, and the idea of position-dependent
CNOT orders is established~\cite{TomitaSvore,KishonyFowler,Strikis2026}\cl{20}.

For the IBM $\code{72,12,6}$ circuit~\cite{Bravyi2024}, which is interleaved,
$\dext=5<6$, so Prop.~3 alone does not suffice. The exact lift test shows that
none of the weight-5 extended logicals lifts, for $R=1,\dots,10$, and a
splitting argument extends this to every $R$: $\dcirc=6$ in both bases, settling
IBM's conjecture for this circuit and fault model\cl{21}. Public data of
Ref.~\cite{Webster2026} already contain an exact value 6 for one basis and $R=6$
of a CZ/SI1000 compilation of an IBM-type circuit; our additions are arbitrary
$R$, both bases, the CNOT circuit, and the method\cl{21}.

The same pipeline---colour the checks by a translation subgroup, prune block
orders by single-colour-class witnesses, enumerate legal states, and prove the
survivors by DFS---gives the theorems in the lower half of Table~\ref{tab:main}
\cl{22} and rigorous negative results: every two-colour LRC of the gross code has
$\dcirc\le10$, and every two- or three-class mod-3 LRC of $\code{90,8,10}$ has
$\dcirc\le8$\cl{23}. To our knowledge these are the first rigorous circuit
distances for BB codes with $d\ge9$\cl{25} (but see the caveat on
Ref.~\cite{Tripier2026} in SM~\ref{sm:novelty}).

% ---------------------------------------------------------------------------
\paragraph{Discussion.}
Three lessons emerge. First, the code is not the circuit: an invariance of the
end-cycle code under a term reordering does not carry over to $\dcirc$, and a
larger end-cycle distance ($\code{72,12,8}$) need not give a larger circuit
distance. Second, reductions must be checked in the form they are used: the premise of
Prop.~1 of Ref.~\cite{ShawTerhal2026} holds for the circuits we treat, but its
per-fault wording is basis dependent, and our frame model describes only the
bulk. Third, exact certification at $n\approx 300$ qubits and $d=12$ is feasible
on a workstation ($10^{10}$--$10^{11}$ search nodes) once symmetry and the
extended-code structure are exploited.

\paragraph{Open questions.}
(i) Does any interleaved single-ancilla circuit for the gross code reach
$\dcirc=12$~\cite{Strikis2026}? Our results exclude the retimings of Fig.~12 in
the families listed above; per-check-free pipelined retimings with $P\ge9$ and
circuits with different CNOT orders remain open. (ii) Is the IBM morphing idling
circuit of Ref.~\cite{Wills2026} the circuit we certify? (iii) Three-round
($\mathbb{Z}_3$) morphing circuits and the non-Abelian codes of
Ref.~\cite{ShawTerhal2025} (App.~C) are not covered by the fold reduction
used here. (iv) Do the proven distances survive correlated faults, leakage, or
CZ-native compilations? (v) Do the distance gains translate into measurable
logical-error-rate gains with practical decoders at accessible $p$? We have not
established this. (vi) Can the remaining gaps---$\code{120,8,12}$ (proven 10),
$\code{90,8,10}$ (non-interleaved $\le 9$), $\code{288,12,18}$ (IBM circuit
$\le 18$, unproven)---be closed?

\begin{acknowledgments}
\todo{Acknowledgements; disclosure of AI-agent involvement in the research and
writing -- PI decision.}
\end{acknowledgments}

% ===========================================================================
\begin{thebibliography}{99}
\bibitem{Bravyi2024} Bravyi, Cross, Gambetta, Maslov, Rall, and Yoder, High-threshold and low-overhead fault-tolerant quantum memory,
Nature \textbf{627}, 778 (2024), arXiv:2308.07915.
% NOTE: first names/initials and most titles are not recorded in the project's
% literature notes; they are omitted rather than reconstructed from memory.
\bibitem{Strikis2026} Strikis, Browne, and Beverland,
arXiv:2603.05481 (v2, 2026). \unverified\ (title, first names).
\bibitem{ShawTerhal2025} M.~H.~Shaw and Terhal, Phys.\ Rev.\ Lett.\
\textbf{134}, 090602 (2025), arXiv:2407.16336. \unverified\ (title).
\bibitem{ShawTerhal2026} M.~H.~Shaw and Terhal, arXiv:2604.09797 (v2, 2026).
\unverified\ (title).
\bibitem{Wills2026} Wills, Beverland, Bishop, Gambetta, Rall, Siddhu, and Cross, arXiv:2605.21898 (v2, 2026), App.~C.
\unverified\ (title, initials).
\bibitem{TourDeGross} Yoder, Schoute, Rall, Pritchett, Gambetta, Cross, Carroll,
and Beverland, Tour de gross: A modular quantum
computer based on bivariate bicycle codes, arXiv:2506.03094 (2025).
\bibitem{Webster2026} Webster, Jacob, and Higgott,
arXiv:2603.22532 (v2, 2026), and the public results table in the repository
\texttt{m-webster/codeDistancePYPI}. \unverified\ (title).
\bibitem{Gidney2021stim} Gidney, Stim: a fast stabilizer circuit simulator,
Quantum \textbf{5}, 497 (2021), arXiv:2103.02202.
\bibitem{Gidney2021honeycomb} Gidney, Newman, Fowler, and Broughton,
A fault-tolerant honeycomb memory, arXiv:2108.10457 (2021).
\bibitem{WangMueller} Wang and Mueller, Coprime bivariate bicycle codes and their
layouts on cold atoms, arXiv:2408.10001. \unverified\ (authors and title taken
from a search-result URL only).
\bibitem{ManesClaes} Manes and Claes, Quantum (2025), arXiv:2308.15520.
\unverified\ (title, volume).
\bibitem{Liang2026} Liang, Gu, Chen, Eisert, and Wang, arXiv:2609.37231 (2026).
\unverified\ (title, initials).
\bibitem{SunCai2026} Sun and Cai, arXiv:2609.39522 (2026). \unverified\ (title,
initials).
\bibitem{TomitaSvore} Tomita and Svore, Phys.\ Rev.\ A \textbf{90},
062320 (2014), arXiv:1404.3747. \unverified\ (title).
\bibitem{KishonyFowler} Kishony and Fowler, arXiv:2603.28852 (2026).
\unverified\ (title, initials).
\bibitem{Ghosh2026} Ghosh, Boche, and Tanggara, arXiv:2610.03429 (2026).
\unverified\ (title, initials).
\bibitem{Tripier2026} Tripier \emph{et al.}, arXiv:2604.19481 (2026).
\unverified\ (title, author list).
\bibitem{Chao2026} Chao, arXiv:2606.12724 (2026). \unverified\ (title).
\bibitem{qldpcchallenge} \texttt{unitaryfoundation/qldpc-challenge}, GitHub
pull request \#724 (2026-08-26).
\end{thebibliography}

% ===========================================================================
\clearpage
\onecolumngrid
\begin{center}
\textbf{\large Supplemental Material}
\end{center}
\setcounter{section}{0}
\renewcommand{\thesection}{S\arabic{section}}
\setcounter{table}{0}
\renewcommand{\thetable}{S\arabic{table}}
\setcounter{equation}{0}
\renewcommand{\theequation}{S\arabic{equation}}

% ---------------------------------------------------------------------------
\section{Fault model and definition of $\dcirc$}\label{sm:noise}
Noise is inserted by a fixed routine (SD6~\cite{Gidney2021honeycomb}): after every
CNOT a two-qubit depolarizing channel; after every reset a flip; before every
measurement a flip; on every qubit not acted on in a time step a single-qubit
depolarizing channel (also in reset and measurement layers). A \emph{fault} is one
Pauli component of one such channel. Two faults with identical detector and
observable signatures are the same DEM mechanism; $\dcirc$ counts mechanisms.
For lower bounds of CSS circuits we use the CSS relaxation: in a $Z$-basis memory
we keep only $Z$-type detectors and the $X$ part of each fault (and vice versa).
This relaxation can only lower the minimum, so it is valid for lower bounds;
upper bounds are always checked on the full circuit with all detectors.
Lower bounds hold for any sub-model (e.g.\ no idle noise). For the morphing
circuits the witnesses of $\dcirc\le\dt$ use idle components after a reset layer
(or, in the first round, initial reset flips); for a no-idle model those upper
bounds rely on the noisy initial data reset and would need re-checking
otherwise\cl{28}.

% ---------------------------------------------------------------------------
\section{Extended code, Prop.~3, and the exhaustive search}\label{sm:dfs}
\emph{Extended code.} For a single-ancilla circuit in which check $c$ applies its
CNOTs in the order $i_1,\dots,i_w$, the residuals are the suffixes
$\{i_l,\dots,i_w\}$, $l=2,\dots,w$. The $Z$-type extended code appends each $Z$-check
residual (as the sum of its data columns) to $(H_X,L_X)$; analogously for $X$.
Zero columns are dropped and identical $(H,L)$ columns merged.
Prop.~3 of Ref.~\cite{Strikis2026}: $\dcirc\ge\dext$ for any single-ancilla CNOT
circuit (interleaved or not); the converse $\dcirc\le\dext$ (their Cor.~1) needs
non-interleaving and is not used for our lower bounds. The key step---each single
fault, propagated to the end of the circuit, has a data image equal to zero or to
exactly one extended-code column---was checked on the evaluated circuits:
$\code{108,8,10}$ ($R=3$, 16\,416 fault components per basis), gross Fig.~12
($R=1,2,3$), IBM $\code{72,12,6}$ ($R=1,\dots,10$) and the four LRC theorems:
0 exceptions, all columns realised\cl{27}.

\emph{DFS.} To decide whether $H_e x=0$, $L_e x\ne0$ has a solution with
$|x|\le W$: a minimum-weight solution is syndrome-minimal; translate it so that it
contains the representative of its lowest orbit; branch on the columns of an odd
syndrome row, forbidding earlier siblings (a partition of the search space); prune
by $\mathrm{depth}+\lceil |s|/w_{\max}\rceil>W$; look up the last column in a hash
table. The symmetry group is computed (translations that are code automorphisms,
map each check's CNOT order to the image check's order, and preserve logical
classes, verified by solving $L\circ T=ML+NH$ with $\mathrm{rank}\,M=k$).
Validation: exact solution counts against unpruned meet-in-the-middle enumeration
on 22 BB72 instances and on $>64$-row direct-sum instances; split-run completeness
(sum of split runs = unsplit run); runs without symmetry for $\code{108,8,10}$
(both sides, $\approx2.5\times10^{10}$ nodes, 0 solutions) and for the four LRC
theorems\cl{19,22}.

\emph{$\code{108,8,10}$ circuit.} Terms indexed $[A_1,A_2,A_3,B_1,B_2,B_3]$ with
$A=(x^3,y,y^2)$, $B=(y^3,x,x^2)$, $(\ell,m)=(9,6)$. $X$ checks: reset at 0, CNOTs at
1--6, measure at 7; term order $[0,1,2,5,3,4]$ on even $y$ rows and
$[2,1,0,5,3,4]$ on odd rows. $Z$ checks: reset at 4, CNOTs at 5,6,7,0,1,2, measure
at 3; term order $[3,4,5,2,1,0]$ on even rows and $[3,4,5,0,1,2]$ on odd rows.
(The $Z$ side applies the $B$ block first; the reversal acts on the $A$ block.)
No data-qubit conflicts, non-interleaved and partitionable.
$\dext^X=\dext^Z=10$ exactly (270 columns per side, 10 orbits of 27;
weight-9 DFS: 0 solutions with and without symmetry; weight 10 attained);
injection-verified weight-10 witness\cl{19}.

\emph{Gross Fig.~12.} 360 merged columns per side (144 of $H_e$-weight 3, 144 of
weight 6, 72 of weight 9); 24 translations ($a+b\equiv0 \bmod 3$) preserve every
check's CNOT order and the logical classes; 15 orbits. Weight-10 DFS: 0 solutions
on the original instance ($\approx2\times10^{10}$ nodes per side), on an
independently built lattice-coordinate instance, and under an order-8 subgroup
(45 orbits, $\approx5\times10^{10}$ nodes per side). Weight-11 solutions realised
as 11 circuit faults for $R=1,2,3$, both bases\cl{10,11}. A run without any
symmetry ($\approx 7\times10^{11}$ nodes) was not performed.

% ---------------------------------------------------------------------------
\section{Exact lift test and arbitrary $R$ for the IBM $\code{72,12,6}$ circuit}\label{sm:lift}
Let $F$ be an undetected logical with $|F|=5=\dext$. By the Table-I property each
fault image is zero or one column; cancelling pairs leaves an extended logical $T$
with $|T|\le|F|$. Since no extended logical of weight $\le4$ exists, all five
images are nonzero and distinct and form a weight-5 extended logical. We list all
of them (36 for the $X$-type, 252 for the $Z$-type extended code; meet-in-the-middle
and DFS agree), group the DEM terms by the image of every circuit location
(checked to list all locations: location counts equal those from \texttt{stim}),
and search for one term per column with zero detector sum and nonzero
observable. Result: 0 for $R=1,\dots,10$, both bases; positive control at weight 6
finds solutions\cl{21}.
\emph{Arbitrary $R$.} (1)~The CNOT windows of consecutive cycles are separated.
(2)~If an undetected logical $F$ has an empty round strictly inside its time span,
it splits into two undetected parts whose observables add, one of them a
lighter undetected logical; for $|F|=5=\dext$ this contradicts Prop.~3, so a
weight-5 logical spans at most 5 consecutive rounds. (3)~The noisy $(R+1)$-round
circuit equals the $R$-round circuit with one block inserted, so such a logical
translates into a circuit with $\le7$ rounds, covered by the exact runs.
Hypotheses (1) and (3) and the detector structure were machine-checked; step (3)
is a written argument\cl{21}.

% ---------------------------------------------------------------------------
\section{Morphing circuits: construction, single-round instance, fold}\label{sm:fold}
\emph{Construction.} Mid-cycle code $\code{288,12,18}$, $\ell=m=12$,
$A=x^3+y^7+y^2$, $B=y^3+x+x^2$; homomorphism $f_x$; $K=\ker f_x$, $|K|=72$. Each
round: reset the 144 previously measured qubits, apply the 3 layers of the
inverse contraction and the 3 layers of the next contraction (each layer a
perfect matching of 144 CNOTs), measure 144 qubits. Every contracting stabiliser
lands on a single same-basis measured qubit; the end-cycle codes have 72+72
weight-9 generators, rank 66+66, $k=12$, qubit degree 5. Our construction is
gate-identical to an independent one for all 3 homomorphisms $\times$ 2 layer
orders $\times$ 2 bases $\times$ 2 starting rounds. Detectors are derived as the
GF(2) kernel of noiseless samples and split into time-local blocks
$D_t$\cl{1}.

\emph{Instance.} The set of DEM terms equals the set of projected fault
signatures ($T=1..5$, both bases). Every round-$r$ fault touches only $D_r$,
$D_{r+1}$. For $T=1..5$ every round's instance is equivalent to the round-1
instance (time shift plus a translation for even $r$; rank conditions and a
bijection on projected faults). The 72 translations by $K$ are automorphisms of
every layer and induce column permutations preserving undetected logicals; 12
orbits of 72\cl{1}.

\emph{Fold condition (basis-free form).} Let $\mu$ map each round-$r$ fault to its
time-reversed partner in round $r-1$ (checked to be a fault for every round-$r$
fault). Let $Z_r$ be the span of detector functionals that vanish on all faults of
rounds $<r$. If $F$ is undetected with last round $r$ and round-$r$ part $E$, then
$E$ is annihilated by $Z_r$, and $F'=(F\setminus E)\oplus\mu(E)$ is undetected with
the same logical value whenever
\begin{equation}
\Psi(E)=\big[H_{\rm all}(E)+H_{\rm all}(\mu E);\,L(E)+L(\mu E)\big]=0
\quad\text{for all $E$ annihilated by $Z_r$,}
\end{equation}
i.e.\ whenever every row of $\Psi$ lies in the row space of $Z_r$ restricted to
round-$r$ faults. A local version replaces $Z_r$ by the functionals in
$\mathrm{span}(D_r,D_{r+1})$ vanishing on round-$(r-1)$ faults (a subspace, hence
stronger) and depends only on two rounds. Both versions hold for $r=2..T$,
$T=2..5$, both bases, including the noisy final round. Iterating the fold
gives, for every $T$, an undetected logical of no larger weight inside one round,
i.e.\ a solution of the single-round instance. The arbitrary-$T$ step therefore
rests on exact checks for $T\le5$ plus periodicity of the local condition and of
the single-round instance (machine-checked hypotheses plus a written argument).
\emph{Remark.} The per-fault statements in the proof of Prop.~1 of
Ref.~\cite{ShawTerhal2026} (``same detectors'', ``same observables'') do not hold
literally in a general detector basis, and a naive block-wise test fails at the
final round; the condition above is the one the fold needs and it holds\cl{3}.

\emph{Search and witness.} Weight-11 DFS, all 12 orbit cases, both bases:
0 solutions on two equivalent instances (sparse basis: $\approx2.8\times10^{10}$
nodes per basis; record basis, different tree: $\approx2.0\times10^{11}$ nodes per
basis). The end-cycle codes have distance exactly 12 (DFS) and minimum stabiliser
weight exactly 9; injecting a weight-12 end-cycle logical as 12 single-qubit
faults after a reset layer gives 0 detectors and flipped observables for every
admissible round and $T$\cl{1,2}.

% ---------------------------------------------------------------------------
\section{The 80-circuit morphing family}\label{sm:morph80}
\emph{Inventory.} 10 (mother code, homomorphism) pairs:
$\code{72,12,6}$ ($f_x,f_y,f_{xy}$), $\code{108,8,10}$ ($f_y$ only),
$\code{144,12,12}$ ($f_x,f_y,f_{xy}$), $\code{288,12,18}$ ($f_x,f_y,f_{xy}$);
$\times$ \{Table IV, Table VIII\} $\times$ 2 orders of $(a_2,a_3)$ $\times$
2 orders of $(b_2,b_3)$. Reconstructed from Tables IV, VIII and IX of
Ref.~\cite{ShawTerhal2025} only, and found gate-identical (80/80) to an
independent implementation; end-cycle distances agree; Table~IX of
Ref.~\cite{ShawTerhal2025} fixes the permutations of the Table~VIII circuits it
lists\cl{5}.

\emph{Verification depth.} The 56 circuits with mother codes
$\code{72}$, $\code{108}$, $\code{144}$ were solved directly on the full $T=4$
circuit (noisy boundaries) in both bases by two independent pipelines (different
detector derivations, DEM-based versus own fault enumeration, different solvers,
plus a tree-free meet-in-the-middle check on five instances): 112/112 agree. For
the 8 sub-$\dt$ circuits the values were recomputed in $96$ settings (SD6 versus
no idle noise; full $T=4$; $T=5$ with noise only in half-cycles 2--4; $T=5$ with
noise only in one interior half-cycle): identical values throughout. Minimum
witnesses consisting only of CNOT faults in one interior half-cycle were injected
into circuits carrying all $X$- and $Z$-type detectors (16/16 pass)\cl{5,8}.
The 24 circuits with mother code $\code{288,12,18}$ were proven by one pipeline
(single-round DFS plus fold check, all flags verified); three of them
($f_x$ rev with $\dt=10$; $f_y$ std and $f_{xy}$ rev with $\dt=12$) were reproduced
independently on single interior half-cycle instances in both bases, and one also
on the full $T=3$ circuit without folding. The arbitrary-$T$ statement for these
circuits relies on the per-circuit fold computation (code and output flags
audited, not independently re-implemented)\cl{9}.

\emph{Wording constraints we observe.} We do not claim that ``the circuit of
Table~II of Ref.~\cite{ShawTerhal2025} has $\dcirc=6$'', because Tables II/IV do not
fix the order of $(b_2,b_3)$. We do not attribute the sub-$\dt$ errors to reset
faults: CNOT-only realisations exist\cl{6,8}. The $\code{72,12,8}$ circuits are
listed in Table~IX of Ref.~\cite{ShawTerhal2025} with bound $\le6$; our result is
tightness\cl{7}.

% ---------------------------------------------------------------------------
\section{Gross code: lift structure, evaluator family, and bulk obstructions}\label{sm:gross12}
\emph{Weight-11 logicals.} Exhaustive enumeration with two different DFS programs
gives 4224 weight-11 extended logicals per side, all in free orbits of the
24-element group (176 orbits); 360 of them consist of ten data columns and one
length-2 residual (``10 data + r2''). An exact lift checker augments the circuit
with one observable per row of $H$ and $L$ on the final data readout, so each DEM
term carries its extended-code image; for each extended logical a frontier
dynamic program over the 11 column groups decides liftability exactly. On Fig.~12
all $4224/4224$ lift for $R=1,2,3$ (independently reproduced; the original run
covers $R=1..6$)\cl{13}.

\emph{The benchmark family ($P\le8$).} Rules: one period $P\in[2,8]$ shared by
$X$ and $Z$ checks; per check one reset, six CNOTs at distinct offsets
$0<o_1<\dots<o_6<o_{\rm meas}<P$ and one measurement; each data qubit in at most
one CNOT per time step mod $P$; determinism checked on the circuit. Weight 6
forces $P=8$, $o_k=k$, $o_{\rm meas}=7$; keeping the Fig.~12 orders leaves only
$s_c\in\mathbb{Z}_8$ per check. A non-SAT backtracking search (forward checking,
pigeonhole per data qubit, no symmetry) finds exactly 4 solutions with one start
fixed and 32 $=8\times4$ without; two independent SAT encodings agree. All four
are colour-uniform, deterministic, non-interleaved; one is Fig.~12. The other
lattice isomorphism of Fig.~12 is related by a code relabelling, so the
statement does not depend on that choice. Weight-11 witnesses were injected
through the benchmark's own verification routine: 72/72 (4 circuits $\times$
$R=1..6$ $\times$ 2 bases) fire no detector and flip an observable, with 11
distinct DEM terms\cl{14,15}. Not covered: other CNOT orders; $R\ge7$ for this
finite-$R$ statement.

\emph{Frame model.} A data $X$ error at time $\theta$ is first seen by $Z$ check $c$
(readout times $T_c+nP$) in instance $H_c=\min\{n:T_c+nP>\theta\}$. Three readouts
of a data qubit give three length-$P$ open intervals; pairwise intersection is a
difference constraint on the $H_c$, and by one-dimensional Helly a common time
exists, at which an SD6 fault location is available. A hook (ancilla $X$ after the
fourth CNOT of an $X$ check) fixes its arrival. Feasibility of the difference
constraints (no negative cycle) implies a genuine weight-11 undetected error in the
bulk; this was checked by 75 stim injections. The converse is not used. The model
agrees with the exact lift checker in the bulk but not at the initial time
boundary: on a $P=16$ circuit, 96 logicals blocked in the model lift at $R=2,4,6$
via round-0 detectors\cl{18}.

\emph{Certificates.} Blocking every ``10 data + r2'' logical requires a
negative cycle for each; CP-SAT (path encoding, and independently an unrolled
Bellman--Ford encoding) and a CEGAR SAT loop (lazy determinism clauses; nogoods
from the potentials of liftable logicals, with 24 translated images) give
INFEASIBLE/UNSAT for the families in Table~\ref{tab:gross}\cl{16}.

\begin{table}[h]
\caption{Bulk obstruction status for retimings of Fig.~12 that keep its per-check
CNOT orders. ``Proven'' means: every legal deterministic circuit of the family has
a bulk weight-11 undetected error, hence $\dcirc\le11$ for sufficiently many rounds.}
\label{tab:gross}
\begin{ruledtabular}
\begin{tabular}{lll}
family & status & evidence\\
\hline
symmetric (24-translation) non-wrapping frame, any $P$ & proven & CP-SAT, compressed to $P=38$, two encodings\\
symmetric pipelined, $P=8,9,10$ & proven & two CP-SAT encodings; exhaustive for $P=8,9$\\
symmetric pipelined, $P=11$ & proven & CP-SAT and CEGAR\\
symmetric pipelined, $P=12$ & proven, single encoding & CP-SAT (unrolled Bellman--Ford)\\
$H$-invariant pipelined, $P=9$: $|H|=12$ (3), 8, 6 (3), 4 ($\langle(0,3),(6,0)\rangle$, $\langle(3,3)\rangle$), 3 & proven & CEGAR (two families re-run)\\
$H$-invariant pipelined, $P=10$: $|H|=12$ $\langle(1,2)\rangle$, 8, 6 $\langle(2,4)\rangle$ & proven & CEGAR\\
per-check-free non-wrapping frame, $P\le16$ & proven, one encoding & CP-SAT (reproduced; independent encoding UNKNOWN)\\
$|H|=4$ $\langle(3,0)\rangle$ pipelined, $P=9$ & open & UNKNOWN\\
symmetric pipelined, $P\ge13$ & open & --\\
all other pipelined families with $P\ge9$ (incl.\ $|H|=2$, per-check-free, unlisted subgroups at $P=10$) & open & per-check-free $P=9$: SAT too slow\\
per-check-free non-wrapping frame, $P>16$ & open & --\\
different CNOT orders & open & --\\
\end{tabular}
\end{ruledtabular}
\end{table}
\cl{16,17}
\emph{Classical control.} Over 2000 SAT-enumerated and $\approx1.1\times10^5$
random-walk legal deterministic per-check $P=9$ circuits, all 720 single-hook
logicals remained liftable in the frame model. This is a classical baseline, not a
proof\cl{17}.

% ---------------------------------------------------------------------------
\section{Other BB codes}\label{sm:other}
Codes: $\code{72,12,6}$ and $\code{108,8,10}$ as in Ref.~\cite{Bravyi2024}
($A=x^3+y+y^2$, $B=y^3+x+x^2$); coprime $\code{126,12,10}$ ($\ell=7$, $m=9$,
$A=1+\pi+\pi^{58}$, $B=1+\pi^{13}+\pi^{41}$, $\pi=xy$; Ref.~\cite{WangMueller},
Table 2); $\code{120,8,12}$ with $(\ell,m)=(20,3)$, $A=1+x^3y^2+x^7y^2$,
$B=1+x^{16}y+x^9y$ (found in our own code screen; its equivalence to the
$\code{120,8,12}$ code of other works is \unverified); $\code{124,10,10}$ with
$(\ell,m)=(31,2)$, $A=1+x^{27}y+x^{17}$, $B=1+x^{30}y+x^{12}y$ (own screen).
Term orders (indices $[A_1,A_2,A_3,B_1,B_2,B_3]$): $\code{72,12,6}$ two-colour:
as for $\code{108,8,10}$; cp126 uniform: $X$ $[1,0,2,4,3,5]$, $Z$ $[4,3,5,1,0,2]$;
$\code{120}$ and $\code{124}$ uniform: $X$ $[0,1,2,3,4,5]$, $Z$ $[3,4,5,0,1,2]$.
Each theorem: DFS at $W=w-1$ with and without symmetry (0 solutions), DFS at $W=w$
finds solutions, injection-verified weight-$w$ witness for $R=1,2,3$ and by the
time-translation argument of SM~\ref{sm:lift} for every $R$\cl{22}.
All 864 proper two-colour LRC states of $\code{72,12,6}$ have $\dext=6$\cl{22}.

\emph{Negative results.} Gross: for all three index-2 colourings, every two-colour
LRC has a verified weight-10 extended logical, hence $\dcirc\le10$ (Cor.~1;
the template is non-interleaved). $\code{90,8,10}$ ($A=x^9+y+y^2$, $B=1+x^2+x^7$,
$(15,3)$): every LRC whose block order is constant on the $h=0$ class and on its
complement, or on each of the three cosets, for all four $\mathbb{Z}_3$
homomorphisms, has $\dext\le8$\cl{23}. For partitionable non-interleaved
circuits (circuits that can be cut in time into alternating $X$-only and
$Z$-only CNOT segments), $\code{90,8,10}$ has $\dcirc\le9$ and
$\code{108,8,10}$ has $\dcirc\le10$ by the single-check argument of Prop.~4 of
Ref.~\cite{Strikis2026}; uniform tilings give $\le8$ ($\code{90}$, any depth) and
$\le9$ ($\code{108}$, depth 8). IBM-type pipelined circuits are pairwise
non-interleaved but not partitionable and are not covered\cl{24}.

% ---------------------------------------------------------------------------
\section{Prior work and novelty caveats}\label{sm:novelty}
\begin{itemize}
\item Gross Fig.~12 $=11$: Ref.~\cite{Strikis2026} v2 still states 11 as a
conjecture and suggests exhaustive methods without executing them; no proof of
any gross-circuit $\dcirc\ge11$ was found in the literature checked (search date
2026-10-07)\cl{12}.
\item IBM $\code{72,12,6}$: partial precedent. The public results table of
Ref.~\cite{Webster2026} contains an exact value 6 (two exact solvers) for a
$Z$-basis, $R=6$, CZ+H/SI1000 circuit of IBM type; it is not stated as a result in
the paper and its gate-by-gate identity with the IBM order was not checked. The
GitHub leaderboard~\cite{qldpcchallenge} lists an upper bound 6/6 for $R=6$\cl{21}.
\item $\code{108,8,10}$: literature values are $\le8$; position- or colour-dependent
CNOT orders have precedents~\cite{TomitaSvore,KishonyFowler,Strikis2026}; the
circuit is an instance of the left-right framework of Ref.~\cite{Strikis2026}.
Rigorous circuit-distance results exist for hypergraph-product
codes~\cite{ManesClaes}, layer codes~\cite{SunCai2026}, and as certificates for
amplified codes~\cite{Liang2026}; an algebraic LP framework gives upper
bounds~\cite{Ghosh2026}\cl{20}.
\item ``First rigorous $\dcirc$ for BB codes with $d\ge9$'': Ref.~\cite{Tripier2026}
lists $\dcirc=9$ (without ``$\le$'') for a $\code{70,6,9}$ BB code; the method is not
stated in the parts read. This must be resolved before the claim is made publicly
\unverified\cl{25}.
\item Morphing: no rigorous lower bound on $\dcirc$ for any morphing circuit was
found (search date 2026-10-08); Ref.~\cite{Chao2026} certifies code distances of
related constructions, not circuit distances\cl{4}. The 80-circuit study (C5--C9)
and the gross obstruction (C13--C17) have not had a dedicated literature check
\unverified.
\end{itemize}

% ---------------------------------------------------------------------------
\section{Reproduction}\label{sm:repro}
A single script re-runs the DFS lower bounds for $\code{108,8,10}$, the four LRC
theorems, IBM $\code{72,12,6}$ and (optionally) gross Fig.~12; a full run passed on
2026-10-07. Each audit report lists its own reproduction commands; witnesses for
negative results are archived. \todo{public repository and DOI -- PI decision.}

\end{document}
